Puisque la personne est choisie au hasard, on est dans une situation d'équiprobabilité.
\begin{enumerate}
  \item 
  
  Il y a 600 touristes, dont 328 séjournent dans un camping, donc \mbox{$P(C)=\dfrac{328}{600}=\dfrac{41}{75}$}
  \item 
  
  $\overline{C}$ est l’événement contraire de $C$, donc \mbox{$P(\overline{C})=1-P(C)=\dfrac{600}{600}-\dfrac{328}{600}=\dfrac{272}{600}=\dfrac{34}{75}$}.
  \item 
  
  On cherche \mbox{$P(E\cap C)$}. Il y a 600 personnes, dont 96 sont à la fois des touristes étrangers et séjournent dans un camping.
  
  Donc \mbox{$P(E\cap C)=\dfrac{96}{600}=\dfrac{4}{25}$}.
  \item

   Il y a 212 personnes qui sont en location, dont 116 français. Donc la probabilité que la personne choisie, sachant qu'elle est en location, soit française, est \mbox{$\dfrac{116}{212}=\dfrac{29}{53}$}.
\end{enumerate}


\needspace{10\baselineskip}
